\chapter{光的偏振}

\problem{5.1}{试确定下面两列光波
\begin{align*}
  \boldsymbol{E}_1 &= A_0\left[\boldsymbol{e}_x\cos(\omega t-kz)+\boldsymbol{e}_y\cos\left(\omega t-kz-\frac{\pi}{2}\right)\right] \\
  \boldsymbol{E}_2 &= A_0\left[\boldsymbol{e}_x\sin(\omega t-kz)+\boldsymbol{e}_y\sin\left(\omega t-kz-\frac{\pi}{2}\right)\right]
\end{align*}
的偏振态。}
\answerarea

\problem{5.2}{为了比较两个被自然光照射的表面的亮度，对其中一个表面直接进行观察，另一个表面通过两块偏振片来观察，两偏振片的透振方向的夹角为 $60^\circ$。若观察到两表面的亮度相同，则两表面实际的亮度比是多少？已知光通过每一块偏振片后损失入射光能量的 $10\%$。}
\answerarea
\clearpage

\problem{5.3}{两个尼科耳棱镜 $\mathrm{N_1}$ 和 $\mathrm{N_2}$ 主截面的夹角为 $60^\circ$，在它们之间放置另一个尼科耳棱镜 $\mathrm{N_3}$，让平行的自然光通过这个系统。假设各尼科耳棱镜对非常光均无吸收，试问 $\mathrm{N_3}$ 和 $\mathrm{N_1}$ 的主截面的夹角为何值时，通过系统的光强最大？设入射光强为 $I_0$，求此时所能通过的最大光强。}
\answerarea

\problem{5.4}{在两个正交的理想偏振片之间有一个偏振片以匀角速度 $\omega$ 绕光的传播方向旋转（见题 5.4 图），若入射的自然光强为 $I_0$，试证明透射光强为
\[I=\frac{I_0}{16}(1-\cos 4\omega t)\]\fig[0.52]{5_4}}
\answerarea
\clearpage

\problem{5.5}{线偏振光入射到折射率为 $1.732$ 的玻璃片上，入射角是 $60^\circ$，入射光的电矢量与入射面成 $30^\circ$ 角，求由分界面上反射的光强占入射光强的百分比。}
\answerarea

\problemfig[0.34]{5.6}{如题 5.6 图所示，自然光从空气入射到水面上，折射光再投向在水中倾斜放置的玻璃片上。若使从水面与玻璃片上的反射光均为线偏振光，求玻璃片和水面的夹角 $\alpha$（$n_{\text{玻}}=1.50$）。}{5_6}
\answerarea
\clearpage

\problem{5.7}{一束理想的线偏振光垂直入射到一方解石晶体上，它的振动面和主截面成 $30^\circ$ 角，两束折射光通过在方解石后面的一个尼科耳棱镜，其主截面与入射光的振动方向成 $50^\circ$ 角。计算两束透射光的相对强度。}
\answerarea

\problem{5.8}{线偏振光垂直入射到一块光轴平行于表面的方解石波片上，光的振动面和波片的主截面成 $30^\circ$ 角。求：（1）透射出来的寻常光和非常光的相对强度为多少？（2）用钠光入射时如要产生 $90^\circ$ 的相位差，波片的厚度至少应为多少？（$\lambda=589\,\mathrm{nm}$）}
\answerarea
\clearpage

\problem{5.9}{有一块平行石英片是沿平行于光轴方向切出的，要把它切成一块黄光的 $1/4$ 波片，问这块石英片应有多厚？石英的 $n_{\mathrm{e}}=1.552$，$n_{\mathrm{o}}=1.543$，$\lambda=589.3\,\mathrm{nm}$。}
\answerarea

\problem{5.10}{（1）线偏振光垂直入射到一个表面和光轴平行的波片，透射出来的光中，原来在波片中的寻常光及非常光产生了大小为 $\pi$ 的相位差，问波片厚度为多少？已知 $n_{\mathrm{o}}=1.544\,2$，$n_{\mathrm{e}}=1.553\,3$，$\lambda=500\,\mathrm{nm}$；（2）问这块波片应怎样放置才能使透射出来的光是线偏振光，而且它的振动面和入射光的振动面成 $90^\circ$ 角？}
\answerarea
\clearpage

\problem{5.11}{线偏振光垂直入射到一块表面平行于光轴的双折射晶片，光的振动面和晶片光轴成 $25^\circ$ 角，问波片中的寻常光和非常光透射出来后的相对强度如何？}
\answerarea

\problem{5.12}{在两个正交尼科耳棱镜 $\mathrm{N_1}$ 和 $\mathrm{N_2}$ 之间垂直插入一块波片，发现 $\mathrm{N_2}$ 后面有光射出。但当 $\mathrm{N_2}$ 绕入射光向顺时针转过 $20^\circ$ 后，$\mathrm{N_2}$ 的视场全暗。此时，把波片也绕入射光顺时针转过 $20^\circ$，$\mathrm{N_2}$ 的视场又亮了。问：（1）这是什么性质的波片；（2）$\mathrm{N_2}$ 要转过多大角度才能使 $\mathrm{N_2}$ 的视场又变为全暗。}
\answerarea
\clearpage

\problem{5.13}{一束圆偏振光，（1）垂直入射到 $1/4$ 波片上，求透射光的偏振状态；（2）垂直入射到 $1/8$ 波片上，求透射光的偏振状态。}
\answerarea

\problem{5.14}{试证明一束左旋圆偏振光和一束右旋圆偏振光，当它们的振幅相等时，合成的光是线偏振光。}
\answerarea
\clearpage

\problem{5.15}{设一方解石晶片沿平行光轴方向切出，其厚度为 $0.034\,3\,\mathrm{mm}$，放在两个正交的尼科耳棱镜间。平行光束经过第一尼科耳棱镜后，垂直地射到波片上，对于钠光（$589.3\,\mathrm{nm}$）而言，晶体的折射率为 $n_{\mathrm{o}}=1.658$，$n_{\mathrm{e}}=1.486$。试问：（1）通过第二尼科耳棱镜后，光束发生的干涉是加强还是减弱？（2）如果两个尼科耳棱镜的主截面是互相平行的，结果又如何？}
\answerarea

\problem{5.16}{单色光通过一尼科耳棱镜 $\mathrm{N_1}$，然后射到杨氏干涉实验装置的两个细缝上，试问：（1）尼科耳棱镜 $\mathrm{N_1}$ 的主截面与图面应成怎样的角度才能使光屏上的干涉图样中的暗条纹为最暗？（2）在上述情况下，在一个细缝前放置一半波片，并将这半波片绕着光线方向继续旋转，问在光屏上的干涉图样有何改变？}
\answerarea
\clearpage

\problem{5.17}{单色平行自然光垂直入射在杨氏双缝上，屏幕上出现一组干涉条纹。已知屏上 $A$、$C$ 两点分别对应零级亮纹和零级暗纹，$B$ 是 $AC$ 的中点，如题 5.17 图所示，试问：（1）若在双缝后放一理想偏振片 $\mathrm{P}$，屏上干涉条纹的位置、宽度会有何变化？$A$、$C$ 两点的光强会有何变化？（2）在一条缝的偏振片后放一片光轴与偏振片透光方向成 $45^\circ$ 的半波片，屏上有无干涉条纹？$A$、$B$、$C$ 各点的情况如何？\fig[0.52]{5_17}}
\answerarea

\problem{5.18}{厚度为 $0.025\,\mathrm{mm}$ 的方解石晶片，其表面平行于光轴，放在两个正交的尼科耳棱镜之间，光轴与两个尼科耳棱镜的主截面各成 $45^\circ$。如果射入第一个尼科耳棱镜的光是波长为 $400\sim 760\,\mathrm{nm}$ 的可见光，问透过第二个尼科耳棱镜的光中，少了哪些波长的光？}
\answerarea
\clearpage

\sproblem{5.19}{把一块切成长方体的 KDP 晶体放在两个正交的偏振片之间组成一个产生泡克耳斯效应的装置。已知电光常量 $\gamma=1.06\times 10^{-11}\,\mathrm{m/V}$，寻常光在该晶体中的折射率 $n_{\mathrm{o}}=1.51$，若入射光波长为 $550\,\mathrm{nm}$，试计算从晶体出射的两束线偏振光相位差为 $\pi$ 时，所需加在晶体上的纵向电压（叫做半波电压）。}
\answerarea

\sproblem{5.20}{将厚度为 $1\,\mathrm{mm}$ 且垂直于光轴切出的石英片放在两个平行的尼科耳棱镜之间，使从第一个尼科耳棱镜出射的光垂直射到石英片上，某一波长的光波经此石英片后，振动面旋转了 $20^\circ$。问石英片厚度为多少时，该波长的光将完全不能通过？}
\answerarea
\clearpage

\sproblem{5.21}{试求使波长为 $509\,\mathrm{nm}$ 的光的振动面旋转 $150^\circ$ 的石英片厚度，已知石英对这种波长的旋光度为 $29.7^\circ\,\mathrm{mm^{-1}}$。}
\answerarea

\sproblem{5.22}{将某种糖配制成浓度不同的 $4$ 种溶液：$100\,\mathrm{cm^3}$ 溶液中分别含有 $30.5\,\mathrm{g}$、$22.76\,\mathrm{g}$、$29.4\,\mathrm{g}$ 和 $17.53\,\mathrm{g}$ 的糖。分别用旋光量糖计测出它们通过每分米溶液转过的角度依次是 $49.5^\circ$、$36.1^\circ$、$30.3^\circ$ 和 $26.8^\circ$。根据这些结果，算出这几种糖的旋光率的平均值是多少？}
\answerarea
\clearpage

\sproblem{5.23}{如题 5.23 图所示装置中，$S$ 为单色光源，置于透镜 $\mathrm{L}$ 的焦点处，$\mathrm{P}$ 为起偏器，$\mathrm{L_1}$ 为此单色光的 $1/4$ 波片，其光轴与偏振器的透振方向成 $\alpha$ 角，$\mathrm{M}$ 为平面反射镜。已知入射到偏振器的光束强度为 $I_0$，试通过分析光束经过各元件后的光振动状态，求出光束返回后的光强 $I$。各元件对光束的损耗可忽略不计。\fig[0.52]{5_23}}
\answerarea

\sproblem{5.24}{一束椭圆偏振光沿 $z$ 方向传播，通过一个线起偏器。当起偏器透振方向沿 $x$ 方向时，透射强度最大，其值为 $1.5I_0$；当透振方向沿 $y$ 方向时，透射强度最小，其值为 $I_0$。（1）当透振方向与 $x$ 轴成 $\theta$ 角时，透射强度为多少？（2）使原来的光束先通过一个 $1/4$ 波片后再通过线起偏器，$1/4$ 波片的轴沿 $x$ 方向。现在发现，当起偏器透光方向与 $x$ 轴成 $30^\circ$ 角时，透过两个元件的光强最大，求光强的最大值，并确定入射光强中非偏振成分占多少？}
\answerarea
\clearpage

\sproblem{5.25}{有下列几个未加标明的光学元件：（1）两个线偏振器；（2）一个 $1/4$ 波片；（3）一个半波片；（4）一个圆偏振器。除了一个光源和一个光屏外不借助其他光学仪器，如何鉴别上述各个元件。}
\answerarea

\sproblem{5.26}{一束汞绿光以 $60^\circ$ 角入射到磷酸二氢钾（KDP）晶体表面，晶体的 $n_{\mathrm{o}}=1.512$，$n_{\mathrm{e}}=1.470$。设光轴与晶面平行，并垂直于入射面，试求晶体中 $\mathrm{o}$ 光与 $\mathrm{e}$ 光的夹角。}
\answerarea
\clearpage

\sproblem{5.27}{通过尼科耳棱镜观察一束椭圆偏振光时，强度随尼科耳棱镜的旋转而改变。当强度为极小值时，在尼科耳棱镜（检偏器）前插入一块 $1/4$ 波片，转动 $1/4$ 波片使它的光轴平行于检偏器的透振方向，再把检偏器沿顺时针方向转动 $20^\circ$ 就完全消光。问（1）该椭圆偏振光是右旋还是左旋的？（2）椭圆长、短轴之比是多少？}
\answerarea

\sproblem{5.28}{研究性课题：请查阅互联网，了解偏振式 3D 技术的新进展，并写成一篇小论文。}
\answerarea
\clearpage
